\begin{abstract}

\begin{center}
\normalsize \textbf{Abstract}\\
\end{center}

\noindent
Slaughterhouse digestate is a residue that normally carries a disposal cost. The aim of this work is to examines whether it can be pyrolysed into biochar, upgraded by steam activation into a sulfur sorbent, and also how the steam route compares with the CO$_2$ activation used in a preceding project work. Biochar was produced from French (and also Algerian) slaughterhouse digestate at 400, 500, and 600\,$^{\circ}$C. These chars are then characterized by bomb calorimetry, CHNS elemental analysis, and N$_2$ adsorption. The French chars were then steam-activated at 700\,$^{\circ}$C and loaded with sulfur from H$_2$S. The French digestate gave the higher solid yield at every temperature, falling from 45.7 to 42.2\,\% against 37.4 to 33.5\,\% for the Algerian material. Ash rose from 22.6\,\% in the digestate to 43.5--50.0\,\% in the chars and held the gross calorific value at 14.1--14.2\,MJ\,kg$^{-1}$, against 25.1--28.2\,MJ\,kg$^{-1}$ on an ash-free basis, so the low heating value reflects mineral dilution rather than incomplete pyrolysis. Ash-free carbon rose from 67.2 to 76.0\,\% and the atomic H/C ratio fell from 0.481 to 0.163. Steam activation at 0.5\,mL\,g$^{-1}$ removed 11--18\,\% of the char mass, the loss falling as the precursor pyrolysis temperature rose. The specific surface area increased from 0.8--32\,m$^2$\,g$^{-1}$ in the pristine chars to 89--146\,m$^2$\,g$^{-1}$ after activation and sulfur loading, with the maximum at the 500\,$^{\circ}$C precursor rather than at the most carbonized char. Every sample took up sulfur, 36--64\,mg\,g$^{-1}$ by weighing and 5.2--6.8\,\% S by CHNS against about 1.5\,\% in the pristine chars, which is two to three times the uptake of the CO$_2$-activated chars of the preceding work, although the feedstocks differed and no surface-area data exist for the CO$_2$ samples. No single pyrolysis temperature suited every purpose: 500\,$^{\circ}$C gave the most developed pore structure for adsorption, 400\,$^{\circ}$C the highest nitrogen content for soil application, and 600\,$^{\circ}$C the lowest H/C ratio for carbon storage, while the ash content rules all three out as a solid fuel. At 146\,m$^2$\,g$^{-1}$ the char stays well below commercial activated carbon, so the realistic case is a low-cost sorbent for removing H$_2$S from the biogas of the same digestion.
\noindent

\end{abstract}


\vspace{0.5em}

\noindent

\vspace{0.8em}